\section{Results: four simple questions}
\label{sec:results}
\paragraph{Status of numerical results.}
Tables~\ref{tab:main}--\ref{tab:floor} preserve the current three-seed measurements as
\emph{preliminary values}. They were produced with held-out-score checkpoint selection and therefore
must be replaced by validation-selected results before submission. They are retained in this working
draft to fix the intended analyses and prevent the narrative from being tailored after rerunning.

\subsection{Q1--Q2: unseen organs are harder and change rankings}
Across all competitive methods, HEST LOOO scores are far below POOLED scores (Table~\ref{tab:main}).
The gap is not explained by a single architecture: it appears for regression, transformer, spatial,
and generative approaches. Model rankings also differ across regimes and datasets. This supports
reporting the two regimes separately rather than treating
pooled patient generalization as evidence of organ transfer.

The size of the gap is itself a benchmark result. For the strongest rows in Table~\ref{tab:main},
POOLED PCC is near $0.79$, whereas LOOO PCC is near $0.46$. A model can therefore look reliable on
new patients from known organs while losing roughly $0.33$ correlation when the organ changes. COAST
makes this deployment gap visible rather than averaging it away.

\begin{table}[t]
\centering\small
\setlength{\tabcolsep}{4pt}
\resizebox{\columnwidth}{!}{%
\begin{tabular}{lcc}
\toprule
Method & HEST LOOO & HEST POOLED\\
\midrule
Ridge (UNI) & 0.4046 & 0.7302\\
Random forest (UNI) & 0.3749 & 0.7158\\
ST-Net~\cite{stnet} & 0.3595 & 0.4169\\
HisToGene~\cite{histogene} & 0.4379 & 0.7766\\
Hist2ST~\cite{hist2st} & 0.4087 & 0.7689\\
Gene-DML-style~\cite{genedml}$^\dagger$ & 0.4527 & 0.7856\\
\midrule
STFlow~\cite{stflow} & 0.4477 & 0.7853\\
STFlow + global morphology & 0.4590 & 0.7937\\
STFlow + local morphology & \textbf{0.4624} & \textbf{0.7950}\\
\bottomrule
\end{tabular}
}
\caption{Preliminary macro Pearson results (three-seed mean). Replace with validation-selected scores
and 95\% confidence intervals before submission. Bold denotes the best reference-model value, not a
claim of universal state of the art. $^\dagger$Feature-matched COAST adaptation; LOOO uses two
completed seeds and POOLED uses three.}
\label{tab:main}
\end{table}

\begin{table}[t]
\centering\small
\setlength{\tabcolsep}{4pt}
\resizebox{\columnwidth}{!}{%
\begin{tabular}{lcc}
\toprule
Method & STImage LOOO & STImage POOLED\\
\midrule
ST-Net & 0.0891 & 0.2944\\
HisToGene & 0.1267 & 0.7292\\
Hist2ST & 0.1421 & 0.7368\\
STFlow & 0.1224 & 0.7115\\
STFlow + global morphology & 0.1283 & 0.7201\\
STFlow + local morphology & 0.1295 & 0.7203\\
\bottomrule
\end{tabular}
}
\caption{Preliminary results on STImage-1K4M. Conditioning improves the fixed STFlow backbone. The
complete model table is provided in the appendix.}
\label{tab:stimage}
\end{table}

\subsection{Q3: Pearson fails when genes are nearly silent}
HEST organs with the lowest PCC also have substantially lower target variance and detection
(Table~\ref{tab:floor}). The counterexample to a sample-size explanation is useful: kidney and
prostate have the most slides yet score poorly, whereas lung and skin have two slides each but high
PCC. The appropriate conclusion is limited: sparse, low-variance targets substantially restrict the
interpretability of PCC; they do not prove that the models themselves are adequate.

\begin{table}[t]
\centering\small
\setlength{\tabcolsep}{4pt}
\begin{tabular}{lrrr}
\toprule
Held-out organ & PCC & Median variance & Spots $>0$\\
\midrule
Skin & 0.78 & 1.47 & 40\%\\
Colorectal & 0.72 & 1.14 & 74\%\\
Lung & 0.68 & 1.30 & 61\%\\
Pancreas & 0.60 & 0.97 & 58\%\\
Breast & 0.56 & 0.88 & 57\%\\
Kidney & 0.15 & 0.071 & 11\%\\
Liver & 0.04 & 0.030 & 6\%\\
Prostate & 0.13 & 0.007 & 1\%\\
\bottomrule
\end{tabular}
\caption{Preliminary per-organ diagnostics for the local reference. Low Pearson correlation co-occurs
with low target variance and sparse detection.}
\label{tab:floor}
\end{table}

At a predeclared 10\% detection threshold, preliminary HEST LOOO PCC rises from 0.4477 to 0.4789
for STFlow and from 0.4624 to 0.4950 for the local reference. Because filtering uses held-out labels,
this score is explicitly secondary and diagnostic; it must not be used for gene selection, checkpoint
selection, or hyperparameter tuning. The final paper will report the full threshold-sensitivity curve
and alternative metrics, not only the favorable threshold.

\subsection{Q4: FiLM improves six representative backbones}
Table~\ref{tab:film_general} applies morphology conditioning to six backbones. The best compatible
conditioner improves all six in both LOOO and POOLED. In LOOO,
the gains are largest for ST-Net ($+0.038$), DeepSpaCE ($+0.015$), and HisToGene ($+0.014$); a strong
residual MLP probe also improves ($+0.006$, $p=0.032$), showing that the effect is not limited to weak
models. Hist2ST improves modestly but does not reach significance. On the recent Gene-DML-style
predictor, global conditioning improves LOOO PCC from .4527 to .4597 and POOLED PCC from .7856 to
.7892. The appendix reports the complete eight-model screening, including neutral and negative
results outside this main comparison.

\begin{table*}[t]
\centering\small
\setlength{\tabcolsep}{4.2pt}
\begin{tabular}{lrrrrrrrr}
\toprule
& \multicolumn{4}{c}{LOOO} & \multicolumn{4}{c}{POOLED}\\
\cmidrule(lr){2-5}\cmidrule(lr){6-9}
Backbone & None & Global & Local & Best $\Delta$ & None & Global & Local & Best $\Delta$\\
\midrule
ST-Net & .3588 & .3493 & \textbf{.3969} & +.0380 & .4169 & .3286 & \textbf{.5886} & +.1717\\
DeepSpaCE & .4521 & .4590 & \textbf{.4672} & +.0151 & .7834 & .7890 & \textbf{.7908} & +.0074\\
HisToGene & .4384 & \textbf{.4526} & .4522 & +.0142 & .7775 & .7830 & \textbf{.7881} & +.0106\\
Hist2ST & .4126 & .4147 & \textbf{.4183} & +.0057 & .7659 & .7697 & \textbf{.7702} & +.0043\\
MLP probe & .4530 & \textbf{.4589} & .4588 & +.0059 & .7871 & .7878 & \textbf{.7930} & +.0059\\
Gene-DML-style$^\dagger$ & .4527 & \textbf{.4597} & .4576 & +.0070 & .7856 & \textbf{.7892} & .7852 & +.0036\\
\bottomrule
\end{tabular}
\caption{Preliminary cross-architecture conditioning intervention on HEST (three-seed means). Bold
marks the best variant within each backbone and regime. Wilcoxon tests on paired fold--seed values for
the best LOOO variant give $p<10^{-4}$ (ST-Net, DeepSpaCE), $p=0.0002$ (HisToGene), $p=0.074$
(Hist2ST), and $p=0.032$ (MLP probe). Gene-DML significance testing awaits the missing seed reruns.
$^\dagger$Feature-matched adaptation of Gene-DML~\cite{genedml}; the LOOO unconditioned and POOLED
local cells contain two seeds, while all other Gene-DML cells contain three. Full screening statistics
are provided in the appendix.
These values require validation-selected confirmation before submission.}
\label{tab:film_general}
\end{table*}

STFlow follows the same context-sensitive pattern: global and local conditioning improve it by
$+0.011$--$0.015$ PCC in LOOO and $+0.008$--$0.010$ in POOLED. Local has the highest mean, but its
difference from global conditioning is not significant. Together, the cross-backbone results support
a narrower finding: repeated morphology injection tends to help models lacking explicit multi-scale
context. The complete screening shows that this is not universal, so we present the relationship as
an observed pattern and a testable design hypothesis rather than a general guarantee.

\subsection{Required robustness analyses}
Before submission we will populate the following preregistered analyses: (i) Spearman, NRMSE, MAE,
and detection AP; (ii) patient/slide bootstrap intervals; (iii) parameter-matched conditioning
controls; (iv) fixed-validation checkpoint selection; (v) macro versus micro aggregation; (vi)
detection-threshold sensitivity; (vii) platform-stratified results; and (viii) qualitative H\&E,
ground-truth, prediction, and error maps selected by a fixed rule. These are acceptance-critical
components of the evaluation claim, not optional supplemental experiments.
